\documentclass{report}
\usepackage{amsmath,amssymb}
\setlength{\parindent}{0mm}
\setlength{\parskip}{1em}
\begin{document}
\begin{center}
\rule{6in}{1pt} \
{\large Jodi Mead \\
{\bf Statistical Tests for Total Variation Regularization Parameter Selection}}

Department of Mathematics \\ Boise State University \\ 1910 University Dr \\ Boise \\ ID 83725-1555
\\
{\tt jmead@boisestate.edu}\end{center}

Total Variation (TV) is an effective method of removing noise in digital
image processing while preserving edges. The choice of scaling or
regularization parameter in the TV process defines the amount of
denoising, with value of zero giving a result equivalent to the input
signal. We explore three algorithms for specifying this parameter based
on the statistics of the signal in the total variation process.
Consequently, TV regularization is viewed as an M-estimator that is
assumed to converge to a well defined limit even if the probability model
is not correctly specified. Using statistical tests to find a
regularization parameter is advantageous for computationally large
problems because a single regularization parameter is found that
satisfies an appropriate statistical test, rather than manually adjusting
the regularization parameter, or iterating it to zero. This is especially
useful for nonlinear problems where an underlying linear problem is
solved iteratively, taking the guesswork out of choosing the
regularization parameter in each iterate.


\end{document}
